\documentclass{article}
\usepackage{geometry}
\geometry{margin=1.5cm, vmargin={0pt,1cm}}
\setlength{\topmargin}{-1cm}
\setlength{\headheight}{12.5pt}
\setlength{\paperheight}{29.7cm}
\setlength{\textheight}{25.3cm}

% useful packages.
\usepackage[UTF8]{ctex}
\usepackage{fancyhdr}
\usepackage{layout}
\usepackage{luacode}

% import common macros
% % % % % % % % % % % % % % % % % % % % % % % % % % % % % % % %
% Common Macros for LaTeX
% % % % % % % % % % % % % % % % % % % % % % % % % % % % % % % %

% Useful packages
\usepackage{amsfonts}
\usepackage{amsmath}
\usepackage{amssymb}
\usepackage{amsthm}
\usepackage{enumerate}
\usepackage{graphicx}
\usepackage{multicol}
\usepackage[ruled,linesnumbered]{algorithm2e}

% Math operators and common symbols
\newcommand{\dif}{\mathrm{d}}
\newcommand{\innerProd}[1]{\left\langle #1 \right\rangle}
\newcommand{\difFrac}[2]{\frac{\dif #1}{\dif #2}}
\newcommand{\pdfFrac}[2]{\frac{\partial #1}{\partial #2}}
\newcommand{\T}{\mathrm{T}}
\newcommand{\norm}[1]{\left\| #1 \right\|}
\newcommand{\abs}[1]{\left| #1 \right|}

% Vector and Matrix shortcuts
\newcommand{\vecTwo}[2]{
  \begin{bmatrix}
    #1 \\
    #2
\end{bmatrix}}
\newcommand{\vecThree}[3]{
  \begin{bmatrix}
    #1 \\
    #2 \\
    #3
\end{bmatrix}}
\newcommand{\vecTwoH}[2]{
  \begin{bmatrix}
    #1 & #2
  \end{bmatrix}
}
\newcommand{\vecThreeH}[3]{
  \begin{bmatrix}
    #1 & #2 & #3
  \end{bmatrix}
}

% Common fields
\newcommand{\R}{\mathbb{R}}
\newcommand{\C}{\mathbb{C}}
\newcommand{\Z}{\mathbb{Z}}
\newcommand{\N}{\mathbb{N}}

% Solutions and proofs
\newenvironment{solution}{
\begin{proof}[解]}{
\end{proof}}

% Utility to get environment variables (requires lualatex)
\newcommand{\getEnv}[2]{\directlua{tex.print(os.getenv("#1") or "#2")}}


\usepackage{listings}
\usepackage{xcolor}

% Configure listings for Python code highlighting
\lstset{
  language=Python,
  basicstyle=\ttfamily\small,
  keywordstyle=\bfseries\color{blue},
  commentstyle=\color{gray},
  stringstyle=\color{red},
  breaklines=true,
  showstringspaces=false,
  tabsize=4,
  frame=single,
  framesep=5pt,
  rulecolor=\color{gray},
  backgroundcolor=\color{white},
  captionpos=b,
  numberstyle=\tiny\color{gray},
  numbers=left,
  stepnumber=5,
}

\lstnewenvironment{plaintext}[1][]{
  \lstset{
    language=none,
    basicstyle=\ttfamily,
    keywordstyle=,
    commentstyle=,
    stringstyle=,
    numbers=none,
    frame=none,
    backgroundcolor=\color{white},
    #1                % 允许用户自行覆盖默认设置
  }
}{}

% Name and uID
\newcommand{\name}{\getEnv{NAME}{NAME}}
\newcommand{\uid}{\getEnv{UID}{UID}}
\newcommand{\course}{数值代数}
\newcommand{\hwNum}{7}

\newcommand{\fl}{\mathrm{fl}}

\begin{document}

\pagestyle{fancy}
\fancyhead{}
\lhead{\name (\uid)}
\chead{\course \#\hwNum}
\rhead{\today}

\section{题目 1}

证明：对任意的 $A \in \mathbb{R}^{m\times n}$，有 $\mathcal{N}(A^T A) = \mathcal{N}(A)$。

\begin{solution}

  我们需要证明两个集合相等，即对任意向量 $x \in \mathbb{R}^n$，有 $x \in \mathcal{N}(A^T
  A)$ 当且仅当 $x \in \mathcal{N}(A)$。

  \textbf{($\Leftarrow$) $\mathcal{N}(A) \subseteq \mathcal{N}(A^T A)$：}

  若 $x \in \mathcal{N}(A)$，则 $Ax = 0$。因此
  \[
    A^T A x = A^T (Ax) = A^T \cdot 0 = 0,
  \]
  故 $x \in \mathcal{N}(A^T A)$。

  \textbf{($\Rightarrow$) $\mathcal{N}(A^T A) \subseteq \mathcal{N}(A)$：}

  若 $x \in \mathcal{N}(A^T A)$，则 $A^T A x = 0$。计算
  \[
    \|Ax\|_2^2 = (Ax)^T(Ax) = x^T A^T A x = x^T \cdot 0 = 0.
  \]
  因为范数非负，$\|Ax\|_2 = 0$ 当且仅当 $Ax = 0$，故 $x \in \mathcal{N}(A)$。

  综上所述，$\mathcal{N}(A^T A) = \mathcal{N}(A)$。

\end{solution}

\section{题目 2}

设
\[
  A =
  \begin{pmatrix} 0 & -1 \\ 3 & 2 \\ 0 & -2
  \end{pmatrix}, \quad
  b =
  \begin{pmatrix} -\frac{1}{3} \\ \frac{5}{3} \\ -\frac{2}{3}
  \end{pmatrix}.
\]

\begin{enumerate}
  \item[(1)] 请写出对应最小二乘问题的正则化方程；
  \item[(2)] 请用平方根法求解上述正则化方程。
\end{enumerate}

\begin{solution}

  \textbf{(1) 正则化方程}

  最小二乘问题 $\min_x \|Ax - b\|_2^2$ 的正则化（法线）方程为
  \[
    A^T A x = A^T b.
  \]

  首先计算 $A^T$：
  \[
    A^T =
    \begin{pmatrix} 0 & 3 & 0 \\ -1 & 2 & -2
    \end{pmatrix}.
  \]

  计算 $A^T A$：
  \[
    A^T A =
    \begin{pmatrix} 0 & 3 & 0 \\ -1 & 2 & -2
    \end{pmatrix}
    \begin{pmatrix} 0 & -1 \\ 3 & 2 \\ 0 & -2
    \end{pmatrix}
    =
    \begin{pmatrix} 9 & 6 \\ 6 & 9
    \end{pmatrix}.
  \]

  计算 $A^T b$：
  \[
    A^T b =
    \begin{pmatrix} 0 & 3 & 0 \\ -1 & 2 & -2
    \end{pmatrix}
    \begin{pmatrix} -\frac{1}{3} \\ \frac{5}{3} \\ -\frac{2}{3}
    \end{pmatrix}
    =
    \begin{pmatrix} 0 + 5 + 0 \\ \frac{1}{3} + \frac{10}{3} + \frac{4}{3}
    \end{pmatrix}
    =
    \begin{pmatrix} 5 \\ 5
    \end{pmatrix}.
  \]

  因此，正则化方程为
  \[
    \begin{pmatrix} 9 & 6 \\ 6 & 9
    \end{pmatrix}
    \begin{pmatrix} x_1 \\ x_2
    \end{pmatrix}
    =
    \begin{pmatrix} 5 \\ 5
    \end{pmatrix}.
  \]

  \textbf{(2) 用平方根法求解}

  对矩阵 $G = A^T A =
  \begin{pmatrix} 9 & 6 \\ 6 & 9
  \end{pmatrix}$ 进行 Cholesky 分解 $G = LL^T$，其中 $L$ 为下三角矩阵。

  计算第一列：$L_{11} = \sqrt{G_{11}} = \sqrt{9} = 3$。

  计算第二行第一列：$L_{21} = \frac{G_{21}}{L_{11}} = \frac{6}{3} = 2$。

  计算第二列对角线：$L_{22} = \sqrt{G_{22} - L_{21}^2} = \sqrt{9 - 4} = \sqrt{5}$。

  因此
  \[
    L =
    \begin{pmatrix} 3 & 0 \\ 2 & \sqrt{5}
    \end{pmatrix}.
  \]

  验证：
  \[
    LL^T =
    \begin{pmatrix} 3 & 0 \\ 2 & \sqrt{5}
    \end{pmatrix}
    \begin{pmatrix} 3 & 2 \\ 0 & \sqrt{5}
    \end{pmatrix}
    =
    \begin{pmatrix} 9 & 6 \\ 6 & 4+5
    \end{pmatrix} =
    \begin{pmatrix} 9 & 6 \\ 6 & 9
    \end{pmatrix}.
  \]

  现在求解 $Gy = b$，即 $LL^T y = A^T b$。

  \textbf{前向替代} ($Lz = A^T b$)：
  \[
    \begin{pmatrix} 3 & 0 \\ 2 & \sqrt{5}
    \end{pmatrix}
    \begin{pmatrix} z_1 \\ z_2
    \end{pmatrix}
    =
    \begin{pmatrix} 5 \\ 5
    \end{pmatrix}.
  \]

  从第一个方程：$3z_1 = 5$，得 $z_1 = \frac{5}{3}$。

  从第二个方程：$2z_1 + \sqrt{5}z_2 = 5$，得 $\sqrt{5}z_2 = 5 - 2 \cdot
  \frac{5}{3} = 5 - \frac{10}{3} = \frac{5}{3}$，
  因此 $z_2 = \frac{5}{3\sqrt{5}} = \frac{5\sqrt{5}}{15} = \frac{\sqrt{5}}{3}$。

  \textbf{后向替代} ($L^T x = z$)：
  \[
    \begin{pmatrix} 3 & 2 \\ 0 & \sqrt{5}
    \end{pmatrix}
    \begin{pmatrix} x_1 \\ x_2
    \end{pmatrix}
    =
    \begin{pmatrix} \frac{5}{3} \\ \frac{\sqrt{5}}{3}
    \end{pmatrix}.
  \]

  从第二个方程：$\sqrt{5}x_2 = \frac{\sqrt{5}}{3}$，得 $x_2 = \frac{1}{3}$。

  从第一个方程：$3x_1 + 2x_2 = \frac{5}{3}$，得 $3x_1 = \frac{5}{3} -
  \frac{2}{3} = 1$，因此 $x_1 = \frac{1}{3}$。

  所以最小二乘问题的解为
  \[
    x =
    \begin{pmatrix} \frac{1}{3} \\ \frac{1}{3}
    \end{pmatrix}.
  \]

\end{solution}

\end{document}
